\section{From the matrix equations to fusion categories}
\label{tail:reconstruction}

The preceding parts prove the linear, quadratic, and cubic equations.
Only the quartic equations remain before reconstruction can be applied.

\begin{lemma}[The quartics follow from the cubics]\label{tail:quartic-lemma}
The matrix $\mathsf A$ satisfies the quartic Evans--Gannon equations.
\end{lemma}

\begin{proof}
Write their quartic residual as
\begin{align*}
 K(g,h,i,k)={}&\sum_{l,m}\mathsf A_{l,m}\mathsf A_{l+g,h}
                  \mathsf A_{h+m,l+i}\mathsf A_{i,k+m}
       -\mathsf A_{h-g,i-g}\ind{k=g}\\
 &\quad+\delta^{-1}\ind{h=0}\mathsf A_{i,k}
       +\delta^{-1}\mathsf A_{g,h}\ind{i=0}.
\end{align*}
Order-three symmetry gives the polynomial identity
\begin{equation}\label{tail:quartic}
 K(g,h,i,k)=\sum_m\mathsf A_{i,k+m}C(-h,h+m,g-h,i)
             +\mathsf A_{i-h,g-h}Q(g-k,i).
\end{equation}
Indeed, replace $\mathsf A_{l+g,h}$ in the double sum by
$\mathsf A_{-h,l+g-h}$. The inner sum becomes
\[
 C(-h,h+m,g-h,i)+\mathsf A_{i-h,g-h}\mathsf A_{m+g,i}
                  -\delta^{-1}\ind{h=0}\ind{m=0}.
\]
Multiply by $\mathsf A_{i,k+m}$ and sum over $m$.
The remaining quadratic sum is
$Q(g-k,i)+\ind{g=k}-\delta^{-1}\ind{i=0}$.
Its correction terms cancel by order-three symmetry:
$\mathsf A_{i-h,g-h}=\mathsf A_{h-g,i-g}$ and
$\mathsf A_{-h,g-h}=\mathsf A_{g,h}$.
This proves \eqref{tail:quartic}, and $Q=C=0$ gives $K=0$.
\end{proof}

\begin{lemma}[Changing the dimension embedding]\label{tail:embedding}
There is a field automorphism $\sigma$ of $\C$ with $\sigma(\delta)=d$.
Applying it entrywise to $\mathsf A$ gives a solution of the same
reconstruction equations with dimension parameter $d$.
\end{lemma}

\begin{proof}
The two roots of $X^2-nX-1$ are $\delta$ and $d$, and its
discriminant $n^2+4$ is not a square for $n\ge3$.
Thus $\Q(\delta)$ has an automorphism interchanging the roots.
Extend it to $\overline{\Q}$, fix a transcendence basis of
$\C/\overline{\Q}$, and extend to the algebraic closure $\C$.
This gives $\sigma$. All reconstruction identities are polynomial
identities over $\Q$ after their nonzero denominators are cleared,
so they are preserved by $\sigma$.
\end{proof}

\begin{proof}[Proof of Theorem~\ref{intro:main}]
The linear identities \eqref{tail:linear}, the quadratics
(Proposition~\ref{foundation:quadratics}), and the quartics
(Lemma~\ref{tail:quartic-lemma}) give exactly
\cite[Eqs.~(4.7)--(4.10)]{EG} with $\omega=1$.
Theorem~2(a) of that paper allows the negative root $\delta$ and
constructs a spherical fusion category with simple objects
$\alpha^a$ and $\rho_a=\alpha^a\rho$, $a\in G$,
the fusion rules \eqref{intro:fusion1}--\eqref{intro:fusion2}, and
categorical dimensions $1$ and $\delta$.

Apply Lemma~\ref{tail:embedding} and reconstruct once more.
The resulting spherical category has categorical dimensions $1$ and $d$,
equal to its Frobenius--Perron dimensions. Its global dimension is
$n(1+d^2)$, also its Frobenius--Perron dimension, so it is pseudo-unitary.
Finally, the categories for different odd $n$ have different ranks $2n$,
and hence are pairwise inequivalent.
\end{proof}

\begin{remark}
Pseudo-unitarity is the conclusion here. Unitarity would additionally
require a Hermitian positive-dimensional reconstruction matrix
\cite[Theorem~2(c)]{EG}; the argument above does not establish that
property. Neither arithmetic information about the coefficient field
nor a second reflected contour calculation is needed for the stated
existence theorem.
\end{remark}
